%% capitoli/introduzione.tex
%% In testa al file, prima di scrivere, incollare l'estratto di
%% output/F4-indice.md relativo a questo capitolo (tesi sostenuta,
%% fonti, lunghezza stimata): serve a controllare gli scostamenti.
%% Capitolo non numerato: niente \label di sezione (non avrebbero numero).

% L'introduzione si scrive per ultima, quando si sa che cosa si e' dimostrato.
% I limiti dell'indagine vanno dichiarati qui e ripresi nelle conclusioni:
% un limite taciuto e' il primo rilievo del relatore.

\chapter*{Introduzione}
\addcontentsline{toc}{chapter}{Introduzione}
\label{cap:intro}

\section*{Oggetto dell'indagine}
\addcontentsline{toc}{section}{Oggetto dell'indagine}


\section*{La domanda di ricerca}
\addcontentsline{toc}{section}{La domanda di ricerca}


\section*{Metodo e materiali}
\addcontentsline{toc}{section}{Metodo e materiali}


\section*{Limiti dell'indagine}
\addcontentsline{toc}{section}{Limiti dell'indagine}


\section*{Piano del lavoro}
\addcontentsline{toc}{section}{Piano del lavoro}


